% GENERATED FILE. Edit canonical lessons, then run npm run generate:books.
% canonical-source-hash: fnv1a64:a4eea97a30d554cc
% canonical-lessons: ES-C48-sintesis-una-historia

\chapter{Synthesis --- A Small Story}
\label{ch:synthesis-story}
\hlchaptermodality{101}

\begin{chapteropening}
\textbf{By the end of this chapter:} \emph{I can narrate four sentences using both past tenses, with a background and a foreground.}

Complete ES-C48-sintesis-una-historia: identify the two layers of the story and say why tuve belongs to the foreground.
\end{chapteropening}

\section[Practice]{Synthesis — a small story}

\label{lesson:ES-C48-sintesis-una-historia}



\begin{quote}
Everything in this book so far has been sentences. A story is different: it needs a background and a foreground at the same time.
\end{quote}

\subsection*{The exchange}
Four sentences. Watch the tenses take turns.

\begin{quote}
\emph{Yo \textbf{estudiaba} español.} \emph{No \textbf{tenía} un libro.} \emph{\textbf{Tuve} un libro en marzo.} \emph{Y \textbf{cuando} \textbf{tuve} el libro, \textbf{hablé} con Ana.}
\end{quote}

Now read only the imperfects: \emph{estudiaba}, \emph{tenía}. Those two sentences tell you \textbf{what things were like} — the situation, with no edges, going on for as long as it needed to.

Now read only the preterites: \emph{tuve}, \emph{hablé}. Those tell you \textbf{what happened} — two events, each with edges, in order.

The story is the two layers together. Neither layer is a story on its own: the imperfects alone are a description that never moves, and the preterites alone are a list of events with no world around them.

\begin{grammarlens}[title={Grammar Lens: what you can now do}]
Notice the third sentence. \emph{Tuve un libro} is not \textquotedblleft{}I had a book\textquotedblright{} — it is the moment you \textbf{got} one, which is why it belongs in the foreground with \emph{hablé} rather than in the background with \emph{tenía}. The same verb appears in both layers of this story, doing two different jobs, and nothing about it is irregular.

That is the whole rung. From here you can narrate: not just report that something happened, but say what the world was like while it did.

It is worth naming the cost. Two tenses learned separately, then three chapters on nothing but the choice between them — because the forms were never the difficulty. English lets you leave this question unanswered; Spanish does not.
\end{grammarlens}

\subsection*{Your turn}
Say these aloud:

\begin{itemize}
\raggedright
  \item the four sentences straight through
  \item only the imperfects, then only the preterites, and hear the two layers
\end{itemize}

\emph{Twice through:}

\begin{tcolorbox}[breakable,colback=teal!4,colframe=teal!35!black,title={Before you move on}]
Which layer says what things were like? (The \textbf{imperfect}.) Which says what happened? (The \textbf{preterite}.) And what can you do now that you could not before? (\textbf{Tell a story} — not just report events, but say what the world was like while they happened.)
\end{tcolorbox}
